\phantomsection % 创建锚点用于超链接
\section*{致谢}
\label{sec:acknowledgements}
\addcontentsline{toc}{section}{致谢} % 添加致谢至目次

逝者如斯夫，不舍昼夜。
时间从不因人的怠惰或勤奋而停滞，硕士研究生的生涯也如约而至地迎来了尾声。
回顾这两年半，我的心中有许多收获和感慨，也有一些遗憾和叹惋。
这一路走来的经历和成长在这刻悉数浮现脑海，是许多人的陪伴和帮助给予我力量，让我走到了今天。
首先我要感谢我的导师叶永林老师。
第一次见到叶老师时，他便与我相谈许多，从生活到学习，从研究方向到学术作风，
初入二所的我心中的茫然也在老师的温言如玉中逐渐消散。
叶老师对我的谆谆教诲不仅在科研学习上，更在人生态度上，
是他教会了我在迷茫困顿时要坚定信念，不畏困难。
此后的学习生活中，在叶老师的言传身教中，
我逐渐体会了一个科研人应当保持勇于探索、坚定求实的作风和态度，
并且我也努力以这样的标准要求自己。

其次我要感谢我的副导师谢春梅老师。
谢老师在我的学习和生活上给予了我莫大的帮助与关心。
硕士期间，从论文的选题调研和开题，到程序的编写与调试，
以及小论文和大论文的撰写与修改，谢老师都一一悉心教导。
每当我学习上有困难和疑惑时，她都耐心为我解疑答惑，指明方向。
谢老师对学术研究一丝不苟和精益求精的态度成为了我往后在学术道路上的榜样。

感谢中山大学孙鹏楠老师在SPH方法上给我的指导和许多帮助，
当我在SPH方法的理论学习、程序编写和数值模拟中遇到困难时，
孙老师总是不辞辛劳地为我答疑解惑。
还要感谢孙老师课题组的吕鸿冠师兄、管祥善师兄、黄晓婷师姐和徐杨师兄在我学习SPH方法时为我提供的帮助。

感谢部门的程小明老师、朱芸芸老师、倪歆韵研究员、
俞俊博士、王泽博士、王琦彬博士、王思雨博士、郑坤博士、
杨吉师姐、汪敬翔师兄、吴凡蕾师姐、魏子阳师兄、屈毫拓师兄等人在我读研期间对我的关心和照顾。感谢大家。

感谢研究生部副主任王丽艳老师和叶荫老师为我读研期间的学习生活提供的许多帮助。
同时也感谢我来到二所认识的2021级同学们，他们在我生病时的照顾、失意时的鼓励、落寞时的陪伴都让我铭记于心。
还要感谢余昕师兄在我学习生活中一直以来的帮助。感谢大家。

感谢各位专家教授在百忙之中对我的论文进行评审并提供宝贵的意见。

最后感谢我挚爱的父母，感谢你们一直以来的无私奉献与默默支持，给了我不断前进的动力，
同时由衷感谢我身边的朋友与亲人们，你们的关怀使我倍受鼓舞，你们的帮助使我渡过难关，不断进步。